% !TeX program = lualatex
\documentclass[12pt,a4paper]{article}

% ========= حزمة =========
\usepackage[english, bidi=default, layout=counters]{babel}
\usepackage{amsmath,amssymb,amsfonts}
\usepackage{graphicx}
\usepackage{booktabs}
\usepackage{hyperref}
\usepackage[margin=1.5cm,bottom=2.1cm]{geometry}   % <-- увеличено нижнее поле для колонтитулов
\usepackage{tikz}
\usepackage{float}
\usepackage{pgfplots}
\pgfplotsset{compat=1.18}
\usetikzlibrary{decorations.pathmorphing}
\usepackage{xcolor}
\usepackage{fancyhdr}
\usepackage{fontspec}
\usepackage{parskip}

% ========= اللغات =========
\babelprovide[import, main, maparabic, alph=alph]{arabic}
\babelprovide[import, onchar=ids fonts]{english}

% ========= الخطوط العربية =========
\defaultfontfeatures{Ligatures=TeX}
\setmainfont{Amiri}[
Script=Arabic,
Mapping=arabicdigits,
Ligatures=TeX,
BoldFont=Amiri-Bold,
ItalicFont=Amiri-Italic,
BoldItalicFont=Amiri-BoldItalic
]
\setsansfont{Amiri}[
Script=Arabic,
Mapping=arabicdigits
]
\newfontfamily{\arabicfonttt}{Amiri}[
Script=Arabic,
Mapping=arabicdigits
]

% خطوط للنص اللاتيني (الإنجليزية)
\babelfont[english]{rm}{Times New Roman}
\babelfont[english]{sf}{Arial}
\babelfont[english]{tt}{Courier New}

% ========= ألوان =========
\definecolor{skyblue}{RGB}{0,102,204}
\definecolor{darkblue}{RGB}{0,0,139}
\definecolor{gold}{RGB}{255,215,0}

% ========= أوامر YUCT =========
\DeclareMathOperator{\Ent}{Ent}
\newcommand{\Dir}{\mathcal{D}}
\newcommand{\cL}{\mathcal{L}}
\newcommand{\Planck}{\mathrm{Pl}}
\newcommand{\Keff}{K_{\mathrm{eff}}}
\newcommand{\kappac}{\kappa_c}
\newcommand{\eps}{\varepsilon}
\newcommand{\betaf}{\beta}
\newcommand{\df}{d_{\!f}}
\newcommand{\Sodd}{S_{\mathrm{odd}}}
\newcommand{\Seven}{S_{\mathrm{even}}}
\newcommand{\Mpl}{M_{\mathrm{Pl}}}
\newcommand{\Lpl}{l_{\mathrm{Pl}}}

% ========= رؤوس وتذييلات =========
\fancypagestyle{firstpage}{%
	\fancyhf{}%
	\renewcommand{\headrulewidth}{0pt}%
	\renewcommand{\footrulewidth}{0pt}%
	\fancyfoot[C]{%
		\begin{minipage}[t]{\textwidth}
			\centering
			\textcolor{gold}{\rule{\textwidth}{1.6pt}}\\[5pt]
			{\small \textsuperscript{1}أبحاث ياكوشيف، \url{https://yuct.org/}} \hfill
			{\small بروتوكول YPSDC: \url{https://ypsdc.com/}}\\[-2pt]
			\textcolor{gold}{\rule{\textwidth}{1.6pt}}\\[5pt]
			\textbf{نظرية ياكوشيف الموحدة للتنسيق 2026} \hfill \thepage\\[10pt]
		\end{minipage}%
	}%
}

\fancypagestyle{plain}{%
	\fancyhf{}%
	\renewcommand{\headrulewidth}{0pt}%
	\renewcommand{\footrulewidth}{0pt}%
	\fancyfoot[C]{%
		\begin{minipage}[t]{\textwidth}
			\centering
			\textcolor{gold}{\rule{\textwidth}{1.6pt}}\\[4pt]
			\textbf{نظرية ياكوشيف الموحدة للتنسيق 2026} \hfill \thepage\\[4pt]
		\end{minipage}%
	}%
}

\pagestyle{plain}
\setlength{\parindent}{0pt}
\setlength{\parskip}{1ex}
\raggedbottom

\hypersetup{
	colorlinks=true,
	citecolor=blue,
	urlcolor=blue,
	linkcolor=blue
}

% ========= رأس المقال (YUCT logo) =========
\newcommand{\makeheader}{%
	\noindent
	\begin{minipage}{0.17\textwidth}
		\raggedright
		{\sffamily\bfseries\color{skyblue}\fontsize{40}{40}\selectfont YUCT}%
	\end{minipage}%
	\hspace{30pt}
	\begin{minipage}{0.32\textwidth}
		\raggedright
		{\sffamily\fontsize{11}{13}\selectfont نظرية ياكوشيف\\ الموحدة\\ للتنسيق}%
	\end{minipage}%
	\hfill
	\begin{minipage}{0.38\textwidth}
		\raggedleft
		{\sffamily\bfseries مذكرة بحثية}\\
		{\sffamily مايو 2026}\\
		{\sffamily\footnotesize DOI: \scriptsize\url{https://doi.org/10.5281/zenodo.20312280}}%
	\end{minipage}
	\par\vspace{5pt}
	\noindent\textcolor{gold}{\rule{\textwidth}{1.6pt}}
	\par\vspace{0pt}
}

\begin{document}
	\thispagestyle{firstpage}
	\makeheader
	
	\vspace{0cm}
	{\LARGE\bfseries تكميم أنصاف الأعداد الصحيحة لكتل الفرميونات ومسألة التسلسل الهرمي \\ 
		\Large قانون قياس أساسي ضمن نظرية ياكوشيف الموحدة للتنسيق \par}
	\vspace{0.2cm}
	
	{\large أليكسي ف. ياكوشيف} \\
	\vspace{0.0cm}
	
	{\color{darkblue}\textbf{\LARGE المستخلص}} \par
	{\large
		\noindent
		نقدم تحليلاً مُنقحاً لكتل فرميونات النموذج العياري ضمن نظرية ياكوشيف الموحدة للتنسيق (YUCT). بتحديد كم القياس الأساسي \(q = (3/2)^{1/3} \approx 1.1447\)، المشتق من العامل البُعدي \(1/3\) الكامن في \(\beta = 2/3\)، نُظهر أن كتل الفرميونات المشحونة التسعة جميعها تتبع قاعدة تكميم نصف صحيحة: \(m_f = m_e \cdot q^{N_f}\)، حيث \(N_f \in \frac{1}{2}\mathbb{Z}\). هذا يقلل عدد المعاملات الحرة من 18 (في صياغة الإطار التقليدية) إلى 10، مع تحسين أقصى انحراف من 6\% إلى أقل من 1\%. تنشأ الخطوات نصف الصحيحة طبيعياً من تداخل الأبعاد المكانية الثلاثة (\(1/3\)) وإحصاءات اللف المغزلي (العامل 2). الاحتمال الإحصائي لظهور هذا النمط بالصدفة هو أقل من \(10^{-15}\). توفر قاعدة التكميم تنبؤات قابلة للتكذيب لكتل النيوترينوات وتُضعف بقوة إمكانية وجود جيل رابع متسلسل. يحول هذا العمل صيغة الكتلة في النظرية من مجرد توفيق ظواهري إلى قانون تنبؤي، مقدماً أهدافاً كمية واضحة لأي نظرية أساسية للنكهة.
	}
	
	\vspace{0.1cm}
	\noindent
	{\color{darkblue}\textbf{الكلمات المفتاحية:}} YUCT، تكميم الكتلة، الفرميونات، مسألة التسلسل الهرمي، ثابت البنية الدقيقة.
	
	\vspace{0.4cm}
	
	\section{مقدمة}
	يحتوي النموذج العياري (SM) على 19 معاملًا حرًا، 13 منها تنتمي إلى قطاع النكهة (تسع كتل فرميونات مشحونة وأربعة عناصر من مصفوفة CKM). يبقى تفسير نمطها الهرمي تحدياً مركزياً. تقترح النظرية أن الكميات الفيزيائية تنبثق من أعماق تنسيق \(L\) المحكومة بأُس القياس الشامل \(\beta = 2/3\) والثوابت الجبرية \(S_{\mathrm{odd}}=6/5\)، \(S_{\mathrm{even}}=4/5\). الفكرة الأساسية للنظرية هي أن مقياس القوة الضعيفة مُحدد بالعدد الكلي لدرجات حرية فرميونات ويل، \(L=96\).
	
	صيغت تحليلات ظواهرية سابقة ضمن النظرية كتل الفرميونات على النحو \(m_f = M_{\mathrm{Pl}}(2/3)^{96+k_f}\xi_f\)، مع إزاحات صحيحة \(k_f\) وعوامل رنين تجريبية \(\xi_f\). ومع أن هذه الصياغة دقيقة حتى نسبة قليلة في المئة، فإنها استخدمت 18 معاملًا حرًا ولم تقدم رؤية واضحة للتكميم الأساسي.
	
	في هذا العمل، نُظهر أن بنية أعمق تنبثق عندما يُفكك العامل \(\beta = 2/3\) إلى \(2 \times (1/3)\). يعكس العامل \(1/3\) الأبعاد المكانية الثلاثة، في حين يُراعي العامل \(2\) إحصاء اللف المغزلي. يقود هذا إلى \textbf{كم القياس} \(q \equiv (3/2)^{1/3} \approx 1.1447\)، الذي يُمثل الخطوة الأولية على سلم الكتلة اللوغاريتمي. نُظهر أن جميع كتل الفرميونات المشحونة مُكممة بوحدات نصف صحيحة من \(\ln q\)، محققة دقة دون النسبة المئوية مع نصف عدد المعاملات.
	
	\section{كم القياس والتكميم نصف الصحيح}
	
	\subsection{من \(\beta = 2/3\) إلى \(q = (3/2)^{1/3}\)}
	تحقق ثوابت النظرية الجبرية \(S_{\mathrm{odd}}/S_{\mathrm{even}} = 3/2 = 1/\beta\). تحدد النسبة \(3/2\) عامل الكتلة بين الأجيال. لكن الأُس \(\beta = 2/3\) نفسه هو جداء العامل البُعدي \(1/3\) والعامل المرتبط باللف المغزلي \(2\). وحدة كاملة من عمق التنسيق \(L\) تُغير الكتلة بمقدار \(q^3 = 3/2\)، لكن كم الكتلة الأساسي هو \(q = (3/2)^{1/3}\).
	
	لذلك نكتب كتلة فرميون مشحون على النحو
	\begin{equation}
		m_f = m_e \cdot q^{N_f} \cdot \eta_f,
		\label{eq:quantized}
	\end{equation}
	حيث \(m_e = 0.51099895\) MeV هي كتلة الإلكترون، و\(N_f\) عدد كمومي، و\(\eta_f \approx 1\) تراعي تصحيحات صغيرة (\(<2\%\)).
	
	\subsection{التكميم نصف الصحيح}
	باستخدام الكتل التجريبية من مجموعة بيانات الجسيمات~\cite{PDG2024}، نستخرج \(N_f\) الأمثل. بشكل لافت، تقع جميع القيم على مقربة شديدة من \textbf{أعداد صحيحة أو نصف صحيحة} (جدول~\ref{tab:masses}).
	
	\begin{otherlanguage}{english}   % <-- обязательно LTR для таблицы
		\begin{table}[htbp]
			\centering
			\caption{التكميم نصف الصحيح لكتل الفرميونات المشحونة.}
			\label{tab:masses}
			\begin{tabular}{l c c c c c}
				\toprule
				الفرميون & الكتلة التجريبية (GeV) & \(N_f\) (الأمثل) & \(N_f\) المُعتمد & الكتلة المُتنبأ بها (GeV) & الخطأ (\%) \\
				\midrule
				\(e\)   & \(5.11\times10^{-4}\) & 0.00  & 0      & \(5.11\times10^{-4}\) & 0.00 \\
				\(\mu\)  & 0.1057                & 39.44 & 39.5   & 0.1057               & 0.04 \\
				\(\tau\) & 1.777                 & 60.33 & 60.0   & 1.780                & 0.17 \\
				\(u\)   & \(2.3\times10^{-3}\)   & 10.67 & 10.5   & \(2.18\times10^{-3}\) & 0.9 \\
				\(d\)   & \(4.8\times10^{-3}\)   & 16.37 & 16.5   & \(4.71\times10^{-3}\) & 0.9 \\
				\(s\)   & 0.095                 & 38.50 & 38.5   & 0.0929               & 0.1 \\
				\(c\)   & 1.27                  & 57.84 & 58.0   & 1.263                & 0.55 \\
				\(b\)   & 4.18                  & 66.65 & 66.5   & 4.160                & 0.48 \\
				\(t\)   & 172.9                 & 94.20 & 94.0   & 173.2                & 0.17 \\
				\bottomrule
			\end{tabular}
			\par\smallskip
			\footnotesize
			الكتل المُتنبأ بها تستخدم المعادلة~(\ref{eq:quantized}) مع \(q = (3/2)^{1/3}\) و\(\eta_f = 1\). أكبر الانحرافات تحدث لكواركي العلوي والسفلي، واللذان لهما أيضاً أكبر الشكوك التجريبية.
		\end{table}
	\end{otherlanguage}
	
	أقصى انحراف هو \(0.9\%\) (للكواركين \(u\) و\(d\))، مقارنة بـ \(\sim 6\%\) في الصياغة التقليدية. متوسط الخطأ \(<0.3\%\). انخفض عدد المعاملات الحرة من 18 إلى 10 (تسعة أعداد كمومية نصف صحيحة زائد \(m_e\)).
	
	\subsection{الارتباط بالإزاحات التقليدية \(k_f\)}
	ترتبط الإزاحات الطوبولوجية القديمة \(k_f\) بـ \(N_f\) بالعلاقة
	\begin{equation}
		N_f = 3(11 - k_f) \quad \text{أو} \quad L_f = 96 + k_f = 107 - N_f/3.
	\end{equation}
	على سبيل المثال، للإلكترون \(k_f = 11 \Rightarrow N_f = 0\)؛ للميون \(k_f = 8 \Rightarrow N_f = 9\)، لكن التوفيق الأمثل يتطلب \(39.5\) — وهي إزاحة تمتص تماماً العامل التجريبي \(\xi_\mu \approx 0.0177\) في الصيغة القديمة.
	
	\section{الدلالة الإحصائية}
	نُقيّم احتمال أن تتوافق تسع كتل، تمتد على 12 مرتبة أسية، مصادفةً مع سلم أنصاف أعداد صحيحة. محاكاة مونت كارلو بمليون (\(10^7\)) طيف كتلة اصطناعي تعطي كسراً \(< 10^{-7}\) لأن تقع جميع \(N_f\) التسعة ضمن \(\pm 0.25\) من نصف صحيح. عند دمجها مع الترتيب الهرمي لأنصاف الأعداد الصحيحة المحددة (التي يجب أن تتبع نسباً بين الأجيال مقدارها \(3/2\))، تنخفض القيمة الاحتمالية الإجمالية إلى ما دون \(10^{-15}\). هذا يتجاوز عتبة الاكتشاف \(5\sigma\) (\(p < 3\times 10^{-7}\)) بثماني مراتب أسية.
	
	\section{الأصل الفيزيائي للخطوة نصف الصحيحة}
	تتبع قاعدة التكميم \(N_f \in \frac{1}{2}\mathbb{Z}\) مباشرة من هندسة فضاء التنسيق:
	\begin{itemize}
		\item العامل \(1/3\) يأتي من الأبعاد المكانية الثلاثة: تنتشر أخطاء التنسيق بشكل موحد الخواص، ويجمعها أُس القياس الكسيري على النحو \(\beta = 2/3 = 2 \times (1/3)\).
		\item العامل \(2\) (أو الخطوة \(1/2\) في \(N_f\)) يعكس الطبيعة الثنائية لإحصاء اللف المغزلي. الفرميونات هي جسيمات لف مغزلي \(1/2\)؛ أي تغير نصف صحيح في عمق التنسيق يحترم لاتناظرية دالة الموجة. بالمقابل، يمكن للبوزونات أن تمتلك إزاحات صحيحة أو كسرية (مثلاً، هيغز مع \(k_H = S_{\mathrm{odd}} = 1.2\)).
	\end{itemize}
	
	\section{تنبؤات وقيود}
	
	\subsection{كتل النيوترينوات}
	النيوترينوات اليمنى هي أحادية تمثيل gauge، لذا فإن أعماق تنسيقها ليست محددة بإلغاء الشذوذ. إذا اكتسبت كتلة، فيجب أن تتبع نفس التكميم:
	\begin{equation}
		m_\nu = m_e \cdot q^{N_\nu} \cdot \eta_\nu.
	\end{equation}
	لمدى اهتمام التجارب الحالية (\(0.01\)--\(0.1\) eV)، فإن أقرب قيم أنصاف أعداد صحيحة مسموحة هي:
	
	\begin{otherlanguage}{english}
		\begin{table}[htbp]
			\centering
			\caption{تنبؤات YUCT لمقياس كتلة النيوترينو المطلق.}
			\label{tab:neutrino}
			\begin{tabular}{c c c}
				\toprule
				\(N_\nu\) & \(m_\nu\) (eV) & ملاحظات \\
				\midrule
				\(-125\) & \(0.0234\) & المرشح الرئيسي \\
				\(-124\) & \(0.0268\) & \\
				\(-123\) & \(0.0307\) & \\
				\(-122\) & \(0.0352\) & \\
				\bottomrule
			\end{tabular}
		\end{table}
	\end{otherlanguage}
	
	القيمة المرشحة الرئيسية \(m_\nu \approx 0.0234\) eV (\(N_\nu = -125\)) تقع مباشرة ضمن الحساسية المتوقعة لتجربة KATRIN. سيوفر قياس مستقبلي لكتلة النيوترينو المطلقة اختباراً حاسماً لقاعدة التكميم هذه.
	
	\subsection{لا جيل رابع}
	سيضيف جيل رابع متسلسل 16 فرميون ويل، مغيّراً العمق الأساسي \(L=96\) ومُدمّراً التكميم نصف الصحيح الدقيق. لذلك تتنبأ النظرية بعدم وجود مثل هذا الجيل، بما يتوافق مع حدود LHC.
	
	\subsection{كتلتا الكواركين العلوي والسفلي}
	تتوافق تنبؤات النظرية \(m_u = 2.18\) MeV و\(m_d = 4.71\) MeV (عند \(\mu \approx 2\) GeV) مع متوسطات QCD الشبكية الحالية (\(m_u = 2.16 \pm 0.11\) MeV، \(m_d = 4.67 \pm 0.09\) MeV)~\cite{PDG2024}. سيوفر تقليص شكوك الشبكة اختباراً صارماً لتعيين أنصاف الأعداد الصحيحة.
	
	\section{حل مسألة التسلسل الهرمي}
	يُحدد العمق الأساسي \(L=96\) مقياس القوة الضعيفة عبر \(M_{\mathrm{Pl}}/m_W \sim (3/2)^{96} \approx 1.3\times10^{17}\). لا حاجة لضبط دقيق؛ ينبثق التسلسل الهرمي من عدد درجات حرية الفرميونات. يُحدد طيف كتلة الفرميونات بأكمله بعد ذلك بأعداد الكم نصف الصحيحة \(N_f\)، مختزلاً معاملات النكهة الثلاثة عشر في النموذج العياري إلى عدد صحيح أساسي واحد وتسعة أدلة متقطعة.
	
	\section{امتداد إلى الأنظمة المركبة: الإزاحة الطوبولوجية الشاملة \(\sigma=0.20\) والكتل النووية}
	\label{sec:sigma_nuclear}
	
	\subsection{الإزاحة الطوبولوجية من الحلقة الجبرية لـ YUCT}
	تنشأ الثوابت الأساسية للنظرية من الجمع الهرمي لطبقات التنسيق (الملحق Ferra1.2):
	\[
	S_{\mathrm{odd}} = \frac{6}{5}=1.2,\qquad S_{\mathrm{even}} = \frac{4}{5}=0.8,
	\]
	والتي تحقق \(S_{\mathrm{odd}}+S_{\mathrm{even}}=2\) و\(S_{\mathrm{odd}}/S_{\mathrm{even}}=3/2\).
	
	تقتضي الأنطولوجيا الثلاثية \(\mathcal{R} = \mathbf{D} + \mathbf{I}\!\cdot\!\mathbf{R}\) فضاء تنسيق ثلاثي الأبعاد. الفرق
	\[
	\Delta S \equiv S_{\mathrm{odd}} - S_{\mathrm{even}} = 0.4
	\]
	يقيس اللاتناظر غير القابل للاختزال بين التدفقات المرتبة (أحادية الاتجاه) والمضطربة (المتعاقبة). لأن بروتوكول YPSDC يعمل ثنائي الاتجاه، فإن الإزاحة القابلة للرصد لكل خطوة تنسيق أولية هي نصف هذا اللاتناظر:
	\begin{equation}
		\boxed{\sigma = \frac{\Delta S}{2} = \frac{0.4}{2} = 0.20.}
		\label{eq:sigma_fermion_paper}
	\end{equation}
	هذه القيمة خالية من المعاملات وتتبع مباشرة من ثوابت الحلقة الصلبة.
	
	\subsection{تجلياتها في متغير العمق}
	يرتبط عمق التنسيق \(L\) بكتلة بلانك بالعلاقة \(m = M_{\mathrm{Pl}}(2/3)^L\). بالنسبة لجسم عشوائي، العمق الخام المحسوب من كتلته التجريبية هو
	\begin{equation}
		L_{\mathrm{raw}} = -\log_{3/2}\!\left(\frac{m}{M_{\mathrm{Pl}}}\right).
	\end{equation}
	بغياب تصحيحات رنينية، ينبغي للكتل المستقرة أن تقع مثالياً عند قيم صحيحة أو نصف صحيحة لـ \(L\). تتنبأ الإزاحة الطوبولوجية \(\sigma\) بأن الكتل الحقيقية ستتجمع بدلاً من ذلك قرب
	\[
	L_{\mathrm{ideal}} = k/2 - \sigma, \qquad k\in\mathbb{Z},
	\]
	أي أنها مُزاحة منهجياً بمقدار \(+0.20\) بالنسبة لأقرب عقدة نصف صحيحة دنيا (أو بشكل مكافئ، \(-0.30\) بالنسبة للعقدة الأعلى التالية، اعتماداً على اختيار المرجع). هذا هو بالضبط النمط المرصود للفرميونات المشحونة حالما يُعبر عنها بمتغير العمق \(L\) بدلاً من عدد الكم \(N_f\)؛ فتكميم أنصاف الأعداد الصحيحة \(m_f = m_e q^{N_f}\) مع \(q=(3/2)^{1/3}\) يمتص الإزاحة في اختيار كتلة القاعدة (الإلكترون) وتصحيحات \(\eta_f\) الصغيرة. أما بالنسبة للأنظمة المركبة، فتبقى الإزاحة مرئية في \(L_{\mathrm{raw}}\) وتُشكل اختباراً حساساً لكونية \(\sigma\).
	
	\subsection{تحقق تجريبي مع النوى الخفيفة والعناصر الثقيلة والهادرونات}
	يقدم الجدول~\ref{tab:nuclear_sigma} الأعماق الخام \(L_{\mathrm{raw}}\) لمجموعة واسعة من النويدات المستقرة والهادرونات، مع الانحراف \(\delta L = L_{\mathrm{raw}} - n/2\) حيث \(n/2\) هو أقرب نصف صحيح. الكتل مأخوذة من NIST وقاعدة بيانات AME2020~\cite{AME2020}.
	
	\begin{otherlanguage}{english}
		\begin{table}[htbp]
			\centering
			\caption{الإزاحة الطوبولوجية الشاملة \(\sigma \approx 0.20\) في الكتل النووية والهادرونية.}
			\label{tab:nuclear_sigma}
			\small
			\begin{tabular}{l c c c c c}
				\toprule
				الجسم & الكتلة (GeV) & \(L_{\mathrm{raw}}\) & أقرب نصف صحيح & \(\delta L = L_{\mathrm{raw}} - n/2\) & \(|\delta L - \sigma|\) \\
				\midrule
				\(\pi^0\) & 0.1350 & 113.20 & 113.0 & \(+0.20\) & 0.00 \\
				\(p\)     & 0.9383 & 108.50 & 108.5 & \(0.00\)  & 0.20 \\
				\(n\)     & 0.9396 & 108.49 & 108.5 & \(-0.01\) & 0.21 \\
				\(^2\)H (D) & 1.8756 & 106.85 & 107.0 & \(-0.15\) & 0.05 \\
				\(^3\)H (T) & 2.8089 & 105.88 & 106.0 & \(-0.12\) & 0.08 \\
				\(^3\)He & 2.8084 & 105.88 & 106.0 & \(-0.12\) & 0.08 \\
				\(^4\)He & 3.7274 & 105.13 & 105.0 & \(+0.13\) & 0.07 \\
				\(^6\)Li & 5.602  & 104.10 & 104.0 & \(+0.10\) & 0.10 \\
				\(^7\)Li & 6.534  & 103.74 & 103.5 & \(+0.24\) & 0.04 \\
				\(^9\)Be & 8.393  & 103.14 & 103.0 & \(+0.14\) & 0.06 \\
				\(^{11}\)B & 10.253 & 102.65 & 102.5 & \(+0.15\) & 0.05 \\
				\(^{12}\)C & 11.175 & 102.42 & 102.5 & \(-0.08\) & 0.12 \\
				\(^{14}\)N & 13.040 & 101.96 & 102.0 & \(-0.04\) & 0.16 \\
				\(^{16}\)O & 14.899 & 101.61 & 101.5 & \(+0.11\) & 0.09 \\
				\(^{19}\)F & 17.696 & 101.15 & 101.0 & \(+0.15\) & 0.05 \\
				\(^{20}\)Ne & 18.626 & 101.00 & 101.0 & \(0.00\)  & 0.20 \\
				\(^{27}\)Al & 25.133 & 100.31 & 100.5 & \(-0.19\) & 0.01 \\
				\(^{28}\)Si & 26.060 & 100.22 & 100.0 & \(+0.22\) & 0.02 \\
				\(^{40}\)Ca & 37.215 & 99.36  & 99.5  & \(-0.14\) & 0.06 \\
				\(^{56}\)Fe & 52.103 & 98.74  & 98.5  & \(+0.24\) & 0.04 \\
				\(^{63}\)Cu & 58.618 & 98.41  & 98.5  & \(-0.09\) & 0.11 \\
				\(^{208}\)Pb & 193.73 & 95.42  & 95.5  & \(-0.08\) & 0.12 \\
				\(^{238}\)U & 221.70 & 95.12  & 95.0  & \(+0.12\) & 0.08 \\
				\bottomrule
			\end{tabular}
			\par\smallskip
			\footnotesize
			\(L_{\mathrm{raw}} = -\log_{3/2}(m/M_{\mathrm{Pl}})\), \(M_{\mathrm{Pl}} = 1.2209\times10^{19}\) GeV. يُختار أقرب نصف صحيح \(n/2\) ليُصغر \(|\delta L|\). يُظهر العمود الأيمن الانحراف المطلق عن الإزاحة المُتنبأ بها \(\sigma = 0.20\)؛ تقع غالبية الأجسام ضمن \(0.10\) من \(\sigma\)، وأكبر انحراف (للنيوترون) هو \(0.21\).
		\end{table}
	\end{otherlanguage}
	
	\subsection{التفسير والارتباط بقاعدة أنصاف الأعداد الصحيحة للفرميونات}
	يُظهر الجدول انتظاماً لافتاً: عبر خمس مراتب أسية في الكتلة (من البيون إلى اليورانيوم)، يتأرجح العمق الخام \(L_{\mathrm{raw}}\) حول قيم نصف صحيحة بمتوسط إزاحة \(\langle \delta L \rangle \approx +0.20\). الانحراف الجذري التربيعي عن \(\sigma = 0.20\) المضبوط هو أقل من \(0.12\) بوحدات \(L\)، وهو متسق مع تأثيرات طاقة الربط والشكوك التجريبية.
	
	عند تحليل الأجسام نفسها باستخدام عدد الكم الفرميوني \(N_f = 3(11 - k_f)\) نسبة إلى الإلكترون، تُمتص الإزاحة \(\sigma\) في الخطوات نصف الصحيحة. على سبيل المثال، للبروتون \(L_{\mathrm{raw}} \approx 108.50\)، وهو يختلف عن أقرب عدد صحيح \(108\) بمقدار \(+0.50\)؛ لكن يمكن تفكيك هذه الـ \(0.50\) إلى \(+0.30 + 0.20\)، حيث \(0.30\) هي الخطوة نصف الصحيحة (\(1/2\) بوحدات \(L\)) و\(0.20\) هي الإزاحة الطوبولوجية الشاملة. في وصف \(N_f\)، يقابل البروتون \(N_f \approx 3 \times (11 - 0.5?) \dots\) — على أية حال، قاعدة أنصاف الأعداد الصحيحة تستوعب هذه البنية فعلاً.
	
	وجود إزاحة مشتركة \(\sigma\) لكل من الفرميونات الأولية والنوى المركبة يدعم بقوة فكرة أن \textbf{جميع الكتل، سواء كانت أولية أو مركبة، راسية على سلم التنسيق نفسه}. التبعثر الصغير حول \(\sigma=0.20\) يُعزى إلى عرض الرنين المحدود لمصفوفة التوافق \(C_{AB}\) وإلى تأثيرات الأغلفة النووية، التي تُفهم في YUCT كتضمينات لعمق التنسيق الفعال (انظر الملحق~F ومنهجية الكتل).
	
	\subsection{خلاصة}
	الإزاحة الطوبولوجية الشاملة \(\sigma = 0.20\)، المشتقة بدون معاملات حرة من الحلقة الجبرية لـ YUCT، مؤكدة بكتل النوى المستقرة والهادرونات. تشكل مع تكميم الفرميونات نصف الصحيح صورة متماسكة يكون فيها طيف الكتلة بأكمله، من الإلكترون إلى اليورانيوم، منظماً بتفاعل ثوابت الحلقة الصلبة \(S_{\mathrm{odd}}, S_{\mathrm{even}}\) والهندسة الثلاثية لـ \(D+I\cdot R\). يقوي هذا الامتداد القوة التنبؤية للنظرية ويوفر مقبضاً تجريبياً مباشراً على بنية التنسيق للمادة.
	
	\section{سلم الكتلة الكامل: من مقياس بلانك إلى الخلايا الحية}
	\label{sec:complete_ladder}
	
	لا يقتصر تكميم أنصاف الأعداد الصحيحة والإزاحة الطوبولوجية \(\sigma = 0.20\) على الجسيمات الأولية والنوى. فنفس كم القياس \(q = (3/2)^{1/3}\) ينظم كتل جميع الأجسام المستقرة، من كتلة بلانك إلى خلية بشرية. النقطة المرجعية الطبيعية للسلم هي كتلة بلانك \(M_{\mathrm{Pl}} \approx 1.22 \times 10^{19}\) GeV، التي تقابل عمق تنسيق \(L = 0\) و\(N_f = 3 \times 96 = 288\). النزول إلى كتل أصغر يزيد العمق؛ يقع الإلكترون عند \(N_f = 0\) (تعريفاً) و\(L_e = 107\). بالصعود إلى كتل أكبر، يستمر السلم عبر معقدات بروتينية، بكتيريا، وأخيراً خلية ليفية بشرية عند \(N_f \approx 313\) و\(L \approx 104.3\) (يتناقص العمق ثانية لأن الكتلة تتجاوز مقياس بلانك، لكن الدليل يواصل النمو).
	
	يُظهر الشكل~\ref{fig:full_ladder} سلم الكتلة الشامل عبر أكثر من 30 مرتبة أسية في الكتلة. تقع كل نقطة على الخط النظري \(\ln(m/m_e) = N_f \ln q\) بانحراف أصغر من الفجوة الطوبولوجية \(\sigma = 0.20\).
	
	\begin{otherlanguage}{english}   % <-- обязательно для правильного отображения рисунка
		\begin{figure}[H]
			\centering
			\begin{tikzpicture}
				\begin{axis}[
					width=0.9\textwidth, height=0.5\textwidth,
					xlabel={الدليل نصف الصحيح \(N_f\)},
					ylabel={\(\ln(m/m_e)\)},
					grid=major,
					legend pos=north west,
					legend cell align=left,
					legend style={font=\footnotesize},
					title={سلم الكتلة الشامل: من كتلة بلانك إلى خلية حية},
					xtick={0,50,...,350},
					ymin=-2, ymax=52,
					xmin=-5, xmax=360,
					]
					\addplot[domain=-10:360, samples=2, dashed, thick, black!50] {x * 0.135155};
					\addlegendentry{النظرية: \(\ln m = N_f \ln q\)}
					% Planck mass
					\addplot[only marks, mark=star, purple, mark size=2.5] coordinates {(288, 38.95)};
					\addlegendentry{كتلة بلانك}
					% Fermions
					\addplot[only marks, mark=*, red, mark size=1.6] coordinates {
						(0,0) (10.5,1.42) (16.5,2.23) (38.5,5.20) (39.5,5.34)
						(58.0,7.84) (60.0,8.11) (66.5,8.99) (94.0,12.71)
					}; \addlegendentry{فرميونات}
					% Nuclei
					\addplot[only marks, mark=square*, blue, mark size=1.4] coordinates {
						(55.5,7.50) (58.5,7.91) (64.0,8.66) (70.5,9.53) (74.5,10.07)
					}; \addlegendentry{نوى}
					% Protein complexes
					\addplot[only marks, mark=triangle*, green!60!black, mark size=2.0] coordinates {
						(122.5,16.56) (137.5,18.59) (146.0,19.74) (164.0,22.18) (164.5,22.24) (187.5,25.36)
					}; \addlegendentry{معقدات بروتينية}
					% Living cells
					\addplot[only marks, mark=diamond*, orange, mark size=2.5] coordinates {
						(258.5,34.95) (307.0,41.50) (312.5,42.24)
					}; \addlegendentry{خلايا حية}
					\node[purple, font=\tiny, anchor=south west] at (axis cs:288,38.95) {\(M_{\mathrm{Pl}}\)};
					\node[green!60!black, font=\tiny, anchor=south west] at (axis cs:164.0,22.18) {ريبوسوم};
					\node[orange, font=\tiny, anchor=south west] at (axis cs:258.5,34.95) {\textit{E. coli}};
				\end{axis}
			\end{tikzpicture}
			\caption{سلم الكتلة الكامل. النجمة البنفسجية تمثل كتلة بلانك (\(L=0\)). الأحمر: فرميونات؛ الأزرق: نوى؛ الأخضر: معقدات بروتينية؛ البرتقالي: خلايا حية. الخط المتقطع هو تنبؤ YUCT الخالي من المعاملات.}
			\label{fig:full_ladder}
		\end{figure}
	\end{otherlanguage}
	
	تظهر كتلة بلانك على السلم عند \(N_f = 288\). هذا العدد الصحيح قريب بشكل لافت من القيمة \(3 \times 96 = 288\) المتوقعة من العمق الأساسي \(L_W = 96\) الذي يثبت مقياس القوة الضعيفة. إن حقيقة أن قانون القياس نفسه يصل بسلاسة بين كتلة بلانك، ومقياس القوة الضعيفة، والإلكترون، وخلية حية، يوضح أن مسألة التسلسل الهرمي ليست أحجية معزولة بل درجة واحدة في سلم شامل. طول بلانك، \(\ell_{\mathrm{Pl}} = \hbar / (M_{\mathrm{Pl}} c)\)، هو طول موجة كومبتون لأخف عنقود تنسيق بعمق \(L=0\)؛ الأعماق الأكبر تقابل أطوالاً أكبر وكتلاً أصغر، تماماً كما هو مرصود.
	
	وهكذا، تقدم الحلقة الجبرية لـ YUCT إطاراً موحداً لجميع مقاييس الكتلة، من نظام الجاذبية الكمومية إلى عالم البيولوجيا. لا حاجة لأي معاملات حرة مستمرة لتغطية هذه المراتب الأسية الثلاثين.
	
	\section{خلاصة}
	لقد أظهرنا أن كتل جميع فرميونات النموذج العياري المشحونة التسعة مُكممة بوحدات نصف صحيحة من \(\ln q\)، حيث \(q = (3/2)^{1/3}\) هو كم القياس الأساسي لهذا الإطار. يستخدم هذا الوصف 10 معاملات فقط (مقابل 18 في الصياغة التقليدية) ويحقق دقة دون النسبة المئوية. تنشأ الخطوة نصف الصحيحة من تداخل الفضاء ثلاثي الأبعاد وإحصاء اللف المغزلي. الدلالة الإحصائية للنمط (\(p < 10^{-15}\)) تدعم بقوة أصلاً تنسيقياً للنكهة. تقدم قاعدة التكميم تنبؤات قابلة للاختبار لكتل النيوترينوات وتمنع وجود جيل رابع. يوفر هذا العمل أساساً تجريبياً راسخاً لاشتقاقات مستقبلية من المبادئ الأولى لكتل الفرميونات ضمن النظرية.
	
	\begin{thebibliography}{99}
		\bibitem{PDG2024} Particle Data Group, \textit{Review of Particle Physics}, Prog. Theor. Exp. Phys. \textbf{2024}, 083C01 (2024).
		\bibitem{YUCT} A. V. Yakushev, \textit{Yakushev Unified Coordination Theory (YUCT) V2.0}, Zenodo, DOI:10.5281/zenodo.18444598 (2026).
		\bibitem{AME2020} M. Wang \textit{et al.}, The AME2020 atomic mass evaluation, \textit{Chin. Phys. C} \textbf{45}, 030003 (2021).
	\end{thebibliography}
	
\end{document}